% Pursuing Stacks (cotes 134-2 to 134-8): archive cote 134-3, site page 52 (PDF page 53), typescript page 52.
% typescript. Draft compared: ps2.tex:2606-2681.
% AI-transcribed by Claude Opus 5.5 (claude-opus-5-5, adaptive thinking, effort high), Message Batches API, batch ps1 (msgbatch_01Rv5hCaNTtUt4JHx1yqoHPL), 2026-11-12; not yet checked by eye.
\scanpage{134-3/52}
\mspage{\hand{52}}
\margin{\unsure{Anal. p.\ 351, Cor.\ 1} \unsure{hay mon Rep?}}
\uline{totally} \uline{aspheric} if it \xout{satisfiesxthe} is \inserted{[``non empty'',]} locally aspheric, and satisfies
moreover the equivalent conditions above. It then turns out:
\uline{Corollary}: Any totally aspheric topos is aspheric.

This follows readily from the definition of asphericity, and the
\hand{\v{}}Cech computation of cohomology of X in terms of a \xout{general} family $(a_i)_{i\hand{\in}I}$
covering e , with the $a_i$ in A. As the mutual products and multiple
products of the $a_i$'s are all aspheric (a fortiori acyclic for any \xout{\illegible{1}}
constant coefficients), the \hand{\v{}}Cech calculation is valid (including of course
in the case of non commutative $H^1$) and yields the desired result.

The condition \xout{o}for a topos to be totally aspheric, in contrast to
\uline{local} asphericity, is an exceedingly strong one. If for instance the
topos is defined by a topological space, which I'll denote by X, then it
is seen immediately that any \inserted{two,} non-empty open subsets of X have a non-empty
intersection, in other words X is an \uline{irreducible space} (i.e. \inserted{not empty and} not the union
of two closed subsets distinct from X). It is well known on the other hand
that any irreducible space is aspheric, and \inserted[typed]{clearly} \xout{that} any non-empty open subset
of an irreducible space is again irreducible. Therefore:

\uline{Corollary} 2: A topological\xout{\unsure{l}} space is totally aspheric (i.e. the correspon-%
ding topos is t.a.) iff it is irreducible. In case X is Hausdorf, this
means also that X is a one-point space.

Let's now translate T 3) in the context of general topoi. We get:
\uline{Proposition} 2 Let X be a \xout{totally aspheric} topos. Then the following two
conditions are equivalent:
a) The ``Lawvere element'' $L_X$ of $\hand{\mathcal{A}}$ = Sh(X) is aspheric \xout{(or, what}
\xout{amounts to the same, aspheric} over the final object\xout{\illegible{1}} $e_X = e$.
b) There exists an object I in $\hand{\mathcal{A}}$ which is aspheric \xout{(or, equivalently,}
\xout{aspheric} over e \deleted{)}, and two sections \inserted{$\delta_0,\delta_1$} of I (over e), \xout{defining}
such that Ker($\hand{\delta}_0,\hand{\delta}_1$)= $\hand{\emptyset}$ , i.e. such that $e_0\hand{\cap}e_1 = \hand{\emptyset}$ , where
$e_0,e_1$ are the subobjects of I defined by $\hand{\delta}_0,\hand{\delta}_1$ .

I recall that the Lawvere element is the one which represents the
functor
\[ F \hand{\longmapsto} \text{set of all subobjects of } F \]
on $\hand{\mathcal{A}}$, it is endowed with two sections over e , corresponding to the
two ``trivial'' subobjects $\hand{\emptyset}$ , e of e , and the kernel of this pair of
sections is clearly $\hand{\emptyset}$ , thus a) $\hand{\Longrightarrow}$ b), \inserted{Conversely,} \xout{which is valid for any ``non}
\xout{empty'' topos (it needs not be totally aspheric)x} \over{B}{b}y the homotopy argument
already used (\inserted{section \unsure{31|3)}} \deleted{p. 39-40}), we readily get that b) implies that \deleted{the constant}
\deleted{map from L to L, factoring into}
